\chapter[Multiple Instances of Aims]{Split large cluster allocation to multiple aims instances}
\label{appendix_multiaims}

Specific tasks need to perform many small jobs rather than a few big ones. The job submission of such tasks on large compute facilities which typically require a few 100 MPI processes can be cumbersome and depend on the given job submitter. Aims allows for running multiple instances of itself and split a global MPI communicator int smaller, independent parts. This allows for a unified approach to supmit one large job and let the splitting happening in aims.   
To use this feature, ou need to compile the multiaims executable by compiling the multiaims target, e.g. by typing make multi.scalapack.mpi.
In addition to the usual requirements you need to specify a C-compiler. 

The multiaims environment needs its jobs distributed in subdirectories labeled 1 to $N$, where $N$ is the total number of jobs. In the main directory one additional control file \textit{multiaims.in} needs to be specified. This file contains two keywords:
You now should have a version of FHI-aims that is capable of tracing the memory
consumption. The usage of this functionality is covered in the documentation of the keywords:

\keydefinition{tasks\_per\_subjob}{multiaims.in}
{
  \noindent
  Usage: \keyword{tasks\_per\_subjob} \texttt{value} \\[1.0ex]  
  Purpose: Specifies the number of processes working on one subtasks. 
  Should be a divisor of total number of processes. 
}

\keydefinition{tasks\_per\_subjob}{multiaims.in}
{
  \noindent
  Usage: \keyword{start\_id} \texttt{value} \\[1.0ex]  
  Purpose: Specifies the first job folder to start with. 
}

